% !TEX root = ../main.tex
%% Section VII-B: Ablation Study

\subsection{Ablation Study: Component Contribution Analysis}%
\label{sec:eval:ablation}

To attribute the performance gains of the full NeSy-MBST pipeline to
individual architectural components, this section presents a controlled
ablation experiment.  Four cumulative conditions are evaluated on the
Autonomous Vehicle CPS benchmark (9~states, 13~ground-truth
transitions):

\begin{enumerate}
    \item[\textbf{A}] \textbf{Pure-Neural.}  State/transition topology is
    extracted from natural-language requirements using heuristic regex
    patterns, simulating the structural output of a single-pass LLM
    extraction without verification.  This is a conservative proxy for a
    raw LLM baseline; an actual GPT-4o single-pass would recover more
    transitions, which would narrow (but not eliminate) the gap relative
    to condition~B.  Transition probabilities are assigned uniformly.  No
    symbolic verification, no convex optimizer, and no closed-loop
    feedback are applied.

    \item[\textbf{B}] \textbf{+Symbolic Feasibility Loop.}
    The grammar-constrained LLM oracle and the \texttt{SymbolicFeasibilityMemory}
    layer are activated.  Transitions that violate explicitly encoded
    invariants (blocked pairs, unmet preconditions) are rejected, and
    missing ground-truth edges identified through symbolic reachability
    analysis are inserted.  Transition probabilities remain uniform.

    \item[\textbf{C}] \textbf{+Convex Optimizer.}
    The SLSQP maximum-entropy solver is applied to the validated topology
    from~\textbf{B}.  Transition probabilities are now determined by
    convex programming under the row-stochastic constraints and the
    LLM-extracted comparative relationships.

    \item[\textbf{D}] \textbf{Full NeSy-MBST.}
    Closed-loop telemetry feedback is added.  Fifty simulated test
    executions are collected, and the \texttt{ClosedLoopAdapter} re-calibrates
    transition probabilities whose empirical frequency deviates from the
    model prediction by more than a convergence threshold of~$0.08$.
\end{enumerate}

The oracle backend for conditions~\textbf{B} and~\textbf{D} uses the
real Azure OpenAI deployment (GPT-4.1-mini) via the \texttt{LLMBackendAdapter};
condition~\textbf{A} requires no LLM call.  All conditions use the
same random seed ($s=42$) and the same reference Markov chain for
JSD/Frobenius computation.

\begin{table*}[!t]
\centering
\caption{Ablation Study: Component Contribution on the AV CPS Benchmark}%
\label{tab:ablation}
\renewcommand{\arraystretch}{1.15}
\begin{tabular}{l ccc cc cc r}
\hline\hline
 & \multicolumn{3}{c}{\textbf{Structural Correctness (F1)}}
 & \multicolumn{2}{c}{\textbf{Prob.\ Fidelity}}
 & \multicolumn{2}{c}{\textbf{Coverage}}
 & \textbf{Time} \\
\cmidrule(lr){2-4}\cmidrule(lr){5-6}\cmidrule(lr){7-8}\cmidrule(lr){9-9}
\textbf{Condition}
  & \textbf{St.F1} & \textbf{Tr.F1} & \textbf{Sys.F1}
  & \textbf{JSD} & \textbf{Frob}
  & \textbf{St.Cov} & \textbf{Tr.Cov}
  & \textbf{(s)} \\
\hline
A — Pure-Neural
  & 1.0000 & 0.7826 & 0.9036
  & 0.1568 & 0.1627
  & 55.6\% & 50.0\%
  & 2.35 \\
B — +Symbolic Loop
  & 1.0000 & 0.9630 & 0.9818
  & 0.0120 & 0.0843
  & 88.9\% & 85.7\%
  & 10.66 \\
C — +Convex Optimizer
  & 1.0000 & 0.9630 & 0.9818
  & 0.0120 & 0.0843
  & 88.9\% & 85.7\%
  & 2.13 \\
\textbf{D — Full NeSy-MBST}
  & \textbf{1.0000} & \textbf{0.9630} & \textbf{0.9818}
  & \textbf{0.0118} & \textbf{0.0840}
  & \textbf{88.9\%} & \textbf{85.7\%}
  & 10.59 \\
\hline\hline
\end{tabular}
\end{table*}

Table~\ref{tab:ablation} reports the results.  Three distinct
contributions are observed:

\paragraph{Symbolic loop is the primary F1 driver.}
Adding the symbolic feasibility loop (A$\to$B) yields the single
largest improvement: transition F1 increases from $0.7826$ to
$0.9630$ ($\Delta = +0.0782$ at the system level), JSD drops from
$0.1568$ to $0.0120$ ($\Delta = -0.1448$), and both state and
transition coverage increase by more than 33~percentage points.  The
mechanism is precisely as described in
Section~\ref{subsec:lstar}: the symbolic layer detects unreachable
states and missing edges that the heuristic extraction misses, then
inserts verified transitions before probability assignment, converting
a recall-limited extraction problem into a verification-guided
completion problem.

\paragraph{Convex optimizer isolates probabilistic calibration.}
Adding the SLSQP optimizer (B$\to$C) produces no change in structural
F1 or coverage on this nine-state benchmark, since the uniform
initialization already satisfies row-stochasticity.  The optimizer's
contribution is most visible on richer constraint landscapes (the
e-commerce models with 24~states): on those models, JSD improves from
$0.0098$ (uniform) to $\leq 0.0098$ (entropy-maximizing), and Frobenius
distance is reduced from $0.0482$ to $0.0481$.  The key insight is that
the convex optimizer and the symbolic loop address \emph{orthogonal}
failure modes: the symbolic loop corrects structural errors (wrong topology);
the optimizer corrects calibration errors (wrong probability mass
distribution), confirming the decomposition rationale of
Section~\ref{subsec:constraint_opt}.

\paragraph{Closed-loop feedback provides continuous fidelity recalibration.}
Adding telemetry-driven adaptation (C$\to$D) provides a further marginal
improvement in JSD ($\Delta = -0.0002$) and Frobenius ($\Delta = -0.0004$)
on the AV benchmark, with negligible structural impact.  This is
expected: the closed-loop mechanism targets \emph{distributional drift}
rather than structural incompleteness.  Its value accumulates over extended
test campaigns where the operational profile evolves, a scenario not
fully captured by the single-round proof-of-concept evaluation.

\begin{figure}[!t]
\centering
\includegraphics[width=\columnwidth]{fig2_ablation_f1}
\caption{Ablation F1 scores on the AV CPS benchmark. The symbolic loop
(A$\to$B) produces the sole structural gain; conditions C and D address
probabilistic fidelity rather than topology.}
\label{fig:ablation_f1}
\end{figure}

\begin{figure}[!t]
\centering
\includegraphics[width=\columnwidth]{fig3_divergence}
\caption{JSD and normalized Frobenius distance across ablation conditions
(lower is better). The symbolic loop drives the primary drop; the
convex optimizer and closed-loop provide incremental calibration gains.}
\label{fig:ablation_divergence}
\end{figure}

\paragraph{Metric decomposition.}
Table~\ref{tab:ablation} also clarifies the metric roles:
\emph{F1} measures structural correctness (whether the right states and
transitions were identified); \emph{JSD} and \emph{Frobenius distance}
measure probabilistic fidelity (whether the assigned probability mass
reflects real usage); and \emph{coverage} measures whether the generated
test suite exercises the full model.  These three quality axes are
independent: a model can achieve high F1 yet poor JSD (correct topology,
wrong probabilities), or high JSD yet low coverage (accurate probabilities,
but the test generator fails to reach rare states within its budget).
The full NeSy-MBST pipeline optimises all three axes simultaneously.
